\documentclass[Orbiter Technical Reference.tex]{subfiles}
\begin{document}

\section{Mesh-to-mesh collision detection and response}
\textbf{EvESpirit}\\
\textbf{June 6, 2026}


\subsection{Introduction}
Orbiter implements a rigid-body collision system for vessel-to-vessel,
vessel-to-base, and vessel-to-scatter (foliage or rocks - part of the
Scatterer system) interactions.
Collision remains optionally toggleable for all three categories within
Launchpad $\to$ Options $\to$ Physics settings.

The \emph{primary} collision detection method is a convex hull system based
on the Gilbert-Johnson-Keerthi (GJK) algorithm~\cite{gjk1988} for
intersection testing and the Expanding Polytope Algorithm (EPA) for
penetration depth computation. Today, this is arguably the most widely used
continuous collision detection algorithm in modern 3D video game physics engines
and robotics simulations.
Every vessel is, by default, treated as a convex collider
(\texttt{bIsConvexCollider~=~true}), meaning its sampled hull point cloud
is wrapped in convex sub-hulls that are tested for overlap via GJK/EPA.
This sub-hull method guarantees a much tighter mesh wrap than a standard convex
evaluation would be capable of, and is likewise commonly used throughout
the wider industry.

A \emph{fallback} point-to-triangle detection path, mathematically
equivalent to the M\"oller-Trumbore ray-triangle intersection
test~\cite{moller1997}, is retained for the case where both vessels in a
pair have opted out of convex collision
(\texttt{bIsConvexCollider~=~false}).

The collision \emph{response} for all paths is handled by a standard
Newtonian rigid-body impulse solver~\cite{ibarra2003}, which computes
normal and tangential impulses and applies them as linear and angular
velocity changes.

To maintain framerates that are relatively equivalent or only mildly
affected within the simulation, the system is split into two phases:
a \emph{broad phase} that cheaply discards non-colliding pairs, and a
\emph{narrow phase} that performs the full geometric intersection test.


\subsection{System architecture}
\label{ssec:architecture}
When a vessel pair survives the broad phase
(Section~\ref{ssec:broadphase}), the narrow-phase path is selected
based on each vessel's \texttt{bIsConvexCollider} flag:

\begin{enumerate}
\item \textbf{Convex-Convex} (both \texttt{bIsConvexCollider = true}).
  Each vessel's hull is decomposed into clustered convex sub-hulls
  (Section~\ref{ssec:hullcloud}).  Every pair of sub-hull slices
  undergoes an AABB overlap check; overlapping pairs are tested with
  the GJK algorithm (Section~\ref{ssec:gjk}) for intersection and,
  upon confirmation, the EPA (Section~\ref{ssec:epa}) to obtain the
  penetration depth, contact normal, and contact point.  This is the
  default and fastest path.

\item \textbf{Convex-Concave} (one vessel convex, the other not).
  The convex vessel's clustered hull slices are tested against the
  concave vessel's raw triangle mesh
  (Section~\ref{ssec:convex_concave}).  Each triangle of the concave
  vessel is transformed to planet-local space and tested against the
  overlapping convex hull slices via GJK/EPA.

\item \textbf{Concave-Concave} (both \texttt{bIsConvexCollider = false}).
  The legacy point-to-triangle path
  (Section~\ref{ssec:narrowphase_legacy}).  Hull points of each vessel
  are tested against the triangles of the other vessel using
  barycentric coordinate containment and M\"oller-Trumbore
  ray-triangle CCD.  This is a symmetric test (hull A vs.\ mesh B
  \emph{and} hull B vs.\ mesh A).
\end{enumerate}


\subsection{Broad phase - bounding-sphere rejection}
\label{ssec:broadphase}
Given two vessels $A$ and $B$ with bounding-sphere radii
$r_A$ and $r_B$, the broad phase operates on their relative position and
velocity vectors:
\begin{equation*}
\vec{p}_{\mathrm{rel}} = \vec{p}_B - \vec{p}_A, \qquad
\vec{v}_{\mathrm{rel}} = \vec{v}_B - \vec{v}_A
\end{equation*}
The combined bounding radius is $r_{\Sigma} = r_A + r_B$.
A pair is immediately discarded if the vessels are separating
\emph{and} already well apart:
\begin{equation*}
\vec{p}_{\mathrm{rel}} \cdot \vec{v}_{\mathrm{rel}} \ge 0
\;\;\text{and}\;\;
|\vec{p}_{\mathrm{rel}}| > r_{\Sigma} + \epsilon
\end{equation*}
where $\epsilon$ is a small safety margin (5\,m).

The 5\,m value was chosen to be larger than typical vessel surface 
irregularities ($\sim$1-2\,m for detailed meshes) while remaining small 
enough to avoid premature engagement at long range.

For approaching vessels ($\vec{p}_{\mathrm{rel}} \cdot \vec{v}_{\mathrm{rel}} < 0$),
the time of closest approach is estimated as
\begin{equation*}
t_{\mathrm{min}} = -\frac{\vec{p}_{\mathrm{rel}} \cdot \vec{v}_{\mathrm{rel}}}
                         {|\vec{v}_{\mathrm{rel}}|^2}
\end{equation*}
If $t_{\mathrm{min}}$ falls within the current simulation time step $\Delta t$,
the predicted closest-approach distance is checked against $r_{\Sigma}$:
\begin{equation*}
|\vec{p}_{\mathrm{rel}} + \vec{v}_{\mathrm{rel}}\,t_{\mathrm{min}}|
> r_{\Sigma} + \epsilon
\;\;\Rightarrow\;\; \text{skip pair}
\end{equation*}
Only pairs that survive both tests proceed to the narrow phase.

When two vessels are close enough to enter the narrow phase, the
simulation time warp is automatically reduced to $1\times$ to ensure
the collision is resolved at full temporal resolution.


\subsection{Hull point cloud construction}
\label{ssec:hullcloud}
Each vessel maintains a dense point cloud sampled from its mesh
geometry.  The cache is built by \texttt{RebuildHullCache()}, which
iterates over all mesh groups and constructs a densely sampled,
spatially clustered representation suitable for both convex hull
wrapping (GJK/EPA) and point-to-triangle testing (fallback).

\subsubsection{Triangle subdivision}
For each triangle in every mesh group, the three vertices (offset by the
mesh displacement) are recursively subdivided by bisecting the longest
edge until no edge exceeds 1.0\,m.  This ensures uniform point density
across the hull surface, preventing gaps in coverage on large triangles
that would otherwise leave regions undetectable.

\subsubsection{BSP clustering}
The dense point cloud for each mesh group is spatially partitioned by a
binary space partition (BSP) that recursively bisects the longest axis
of the point set's bounding box until each cluster fits within a
$5.0 \times 5.0 \times 5.0$\,m box.  An overlap margin of 1.1\,m is
applied at each split to ensure that hull points near the bisection
plane appear in \emph{both} child clusters, sealing potential gaps
between adjacent convex sub-hulls.

\subsubsection{Deduplication}
Within each cluster, points are sorted lexicographically by
$(x, y, z)$ and deduplicated with a tolerance of $10^{-4}$\,m.
This eliminates the redundant vertices introduced by the subdivision
of shared triangle edges.

\subsubsection{Slice assembly}
The clusters are flattened into a single \texttt{m\_hullCacheStatic}
array, sorted by mesh index, group index, and cluster index.  A
corresponding \texttt{HullGroupSlice} structure records the start index,
point count, and per-frame AABB for each cluster.  Each slice also
carries an \texttt{isDockClearZone} flag used for docking exclusion
(Section~\ref{ssec:dock_exclusion}).

\subsubsection{Touchdown point injection}
After the mesh-derived hull cache is transformed to planet-local
coordinates (\texttt{UpdateHullCacheP}), the vessel's touchdown vertices
are appended as an additional slice.  This ensures that deployed landing
gear - whose geometry is typically driven by animation rather than
static mesh data - physically interacts with base and scatter
colliders.

\subsubsection{Animation support}
When the vessel has active animations,
\texttt{ApplyAnimationToHullCache()} replays every animation component's
transform (rotation, translation, scale) onto the cached hull vertices
before the planet-local transform.  This keeps the collision hull
synchronised with visual animations such as gear deployment, cargo bay
doors, and solar panel articulation.


\subsection{Primary: GJK/EPA convex hull collision}
\label{ssec:gjk_epa}

\subsubsection{GJK intersection test}
\label{ssec:gjk}
The Gilbert-Johnson-Keerthi (GJK) algorithm~\cite{gjk1988} determines
whether two convex point sets overlap by iteratively building a simplex
in the \emph{Minkowski difference} $A \ominus B = \{\vec{a} - \vec{b}
\mid \vec{a} \in A,\, \vec{b} \in B\}$.  Two shapes intersect if and
only if the origin lies inside their Minkowski difference.

\paragraph{Support function.}
The core primitive is the \emph{support function}, which returns the
point in a convex set farthest along a given direction:
\begin{equation*}
S_A(\vec{d}) = \arg\max_{\vec{a}\in A}\; \vec{a}\cdot\vec{d}
\end{equation*}
A Minkowski-difference support point is then
\begin{equation*}
S(\vec{d}) = S_A(\vec{d}) - S_B(-\vec{d})
\end{equation*}
The implementation stores both the difference point $\vec{p} = S(\vec{d})$
and the original witness points $\vec{p}_A = S_A(\vec{d})$,
$\vec{p}_B = S_B(-\vec{d})$ for later contact-point reconstruction
in the EPA phase.

\paragraph{Simplex evolution.}
Starting from an initial support point, the algorithm iteratively
adds new support points in the direction of the origin and evolves the
simplex through four stages: point $\to$ line $\to$ triangle $\to$
tetrahedron.  At each stage, the simplex is reduced to the
lowest-dimensional face closest to the origin, and a new search
direction is computed.

\begin{itemize}
\item If the tetrahedron encloses the origin, the shapes \emph{intersect}
  and the simplex is passed to the EPA for depth computation.
\item If a new support point does not pass the origin
  ($S(\vec{d}) \cdot \vec{d} \le 0$), the shapes are \emph{separated}.
\item The algorithm is capped at 64 iterations to prevent degenerate
  infinite loops.
\end{itemize}

\subsubsection{EPA penetration depth}
\label{ssec:epa}
When GJK confirms intersection, the Expanding Polytope Algorithm
determines the minimum penetration depth and contact normal.  Starting
from the GJK simplex (a tetrahedron enclosing the origin), the EPA
maintains a convex polytope whose faces are tracked as a list of
triangular facets.

\begin{enumerate}
\item The closest face to the origin is identified.  Its outward
  normal $\hat{\vec{n}}$ and signed distance
  $d = \hat{\vec{n}} \cdot \vec{p}_{\mathrm{face}}$ define the
  candidate penetration vector.
\item A new Minkowski support point is obtained along $\hat{\vec{n}}$.
\item If the new support point does not extend the polytope beyond
  the closest face by more than 0.01\,m (convergence threshold), the
  penetration depth $d$ and normal $\hat{\vec{n}}$ are final.
\item Otherwise, all faces visible from the new support point are
  removed, their horizon edges are identified, and new faces are
  constructed from the horizon edges to the support point.
\item The process repeats for up to 64 iterations.
\end{enumerate}

\paragraph{Contact point reconstruction.}
The contact point on shape~$A$ is computed via barycentric interpolation
of the witness points ($\vec{p}_A$) stored in the three vertices of the
closest polytope face:
\begin{equation*}
\vec{c}_A = u\,\vec{p}_{A,0} + v\,\vec{p}_{A,1} + w\,\vec{p}_{A,2}
\end{equation*}
where $(u, v, w)$ are the barycentric coordinates of the origin's
projection onto the closest face.  This point is initially in
planet-local space and is subsequently transformed to vessel-local
coordinates for impulse computation.


\subsubsection{Convex-Convex collision path}
When both vessels are convex colliders, the algorithm proceeds as follows:

\begin{enumerate}
\item Both vessels' hull caches are transformed to planet-local space.
\item For every pair of hull slices $(s_A, s_B)$ from the two vessels,
  an AABB overlap test is performed.  Non-overlapping pairs are skipped.
\item If a docking approach is detected
  (Section~\ref{ssec:dock_exclusion}), slices marked as dock clear
  zones are excluded.
\item Each overlapping pair is tested with \texttt{GJK\_Intersect}.
  On intersection, \texttt{EPA\_Penetration} computes the depth, normal,
  and contact point.
\item The deepest penetration across all slice pairs is recorded as the
  collision result.
\end{enumerate}


\subsubsection{Convex-Concave collision path}
\label{ssec:convex_concave}
When exactly one vessel is a convex collider, the concave vessel's raw
triangle mesh is used directly.  For each triangle in every mesh group
of the concave vessel:

\begin{enumerate}
\item The triangle's three vertices are transformed to planet-local
  space via the concave vessel's rotation matrix and position offset.
\item The triangle's AABB is checked against the convex vessel's global
  hull AABB; non-overlapping triangles are skipped.
\item For each hull slice of the convex vessel, a per-slice AABB check
  is performed.
\item On overlap, GJK tests the convex hull slice against the three
  triangle vertices (treated as a degenerate convex shape).  On
  intersection, EPA computes the penetration depth and contact point.
\item Docking exclusion zones are respected: if a docking approach is
  active and the convex vessel's slice is a dock clear zone, triangles
  near the concave vessel's dock ports (within 3\,m) are skipped.
\end{enumerate}

The deepest penetration across all triangle-slice pairs is recorded.
The collision normal is oriented consistently from vessel~$B$ toward
vessel~$A$ regardless of which vessel is convex.


\subsection{Fallback: Point-to-triangle intersection (concave-concave)}
\label{ssec:narrowphase_legacy}
When both vessels have \texttt{bIsConvexCollider = false}, the system
falls back to a direct point-to-triangle test using the same hull point
cloud (Section~\ref{ssec:hullcloud}) but without convex wrapping.
This path is symmetric: hull points of $A$ are tested against meshes
of~$B$, and vice versa.

\subsubsection{Coordinate transformation}
The cached hull vertices from vessel~$A$ are transformed into the local
mesh space of vessel~$B$ using the composite rotation matrix
$\mat{R} = \mat{R}_{\mathrm{planet}}^{-1}\,\mat{R}_{\mathrm{vessel}}$.
The inverse transformation is also applied (hull of $B$ tested against
meshes of~$A$) for symmetric detection.

\subsubsection{AABB culling}
Before any per-triangle test, an Axis-Aligned Bounding Box (AABB) overlap
check is performed at two levels:

\begin{enumerate}
\item \textbf{Mesh group AABB.} The AABB of each mesh group in the target
      vessel is compared against the AABB of the incoming hull point cloud.
      Non-overlapping groups are skipped entirely.
\item \textbf{Per-triangle AABB.} Each triangle's bounding box is tested
against individual hull points. A small tolerance of 0.1\,m is added
to both AABBs to account for numerical precision loss in the 
coordinate transformations between vessel-local and mesh-local spaces.
Without this tolerance, edge-case contacts (vertices exactly on AABB
faces etc.) would be incorrectly culled due to floating-point rounding.
\end{enumerate}

\subsubsection{Point-in-triangle test via barycentric coordinates}
\label{sssec:barycentric}
For each triangle with vertices $\vec{p}_0$, $\vec{p}_1$, $\vec{p}_2$
and a candidate hull point $\vec{q}$, define the edge vectors
\begin{equation*}
\vec{e}_1 = \vec{p}_1 - \vec{p}_0, \qquad
\vec{e}_2 = \vec{p}_2 - \vec{p}_0
\end{equation*}
and the triangle face normal
\begin{equation*}
\hat{\vec{n}} = \frac{\vec{e}_1 \times \vec{e}_2}{|\vec{e}_1 \times \vec{e}_2|}
\end{equation*}
The signed distance from the point to the triangle plane is
\begin{equation}
d = (\vec{q} - \vec{p}_0) \cdot \hat{\vec{n}}
\label{eq:signed_dist}
\end{equation}
The signed distance determines the detection path:
\begin{itemize}
\item $d \in [-2.0,\, 0.05]$\,m: the point is within the contact window
  and proceeds to the barycentric test.
\item $d > 0.05$\,m: the point is beyond the triangle surface.  If the
  relative velocity is approaching ($\vec{v}_{\mathrm{rel}} \cdot
  \hat{\vec{n}} < -0.01$), the point may have tunneled through during
  this frame, and a CCD ray-triangle sweep test is performed
  (Section~\ref{sssec:ccd}).
\item $d < -2.0$\,m: the point is too deep to be a valid contact event
  and is rejected.
\end{itemize}

The barycentric coordinates $(u, v, w)$ of the projected point are computed
via the Gram matrix of the edge vectors:
\begin{equation}
\begin{aligned}
d_{00} &= \vec{e}_1 \cdot \vec{e}_1, \quad
d_{01} = \vec{e}_1 \cdot \vec{e}_2, \quad
d_{11} = \vec{e}_2 \cdot \vec{e}_2 \\[4pt]
d_{20} &= (\vec{q} - \vec{p}_0) \cdot \vec{e}_1, \quad
d_{21} = (\vec{q} - \vec{p}_0) \cdot \vec{e}_2 \\[4pt]
D &= d_{00}\,d_{11} - d_{01}^{\,2}
\end{aligned}
\label{eq:gram}
\end{equation}
The barycentric weights are then
\begin{equation}
v = \frac{d_{11}\,d_{20} - d_{01}\,d_{21}}{D}, \qquad
w = \frac{d_{00}\,d_{21} - d_{01}\,d_{20}}{D}, \qquad
u = 1 - v - w
\label{eq:bary}
\end{equation}
A hit is registered when $u \ge -0.05$, $v \ge -0.05$, $w \ge -0.05$.
The small negative tolerance ($-0.05 \approx 5$\% of triangle edge length for
typical meshes) accommodates floating-point imprecision and slight
coordinate misalignment from the dual-direction testing.
The system tracks the deepest penetration depth $-d$ across all tested
point-triangle pairs and records the corresponding surface normal and
contact location.


\subsubsection{Continuous collision detection (CCD) via ray sweep}
\label{sssec:ccd}
When a hull point's signed distance $d > 0.05$\,m, it has passed beyond
the triangle surface.  For thin geometry (e.g.\ wing cross-sections), this
can occur when a hull point tunnels completely through a surface between
frames.  To detect such through-penetrations, a M\"oller-Trumbore
ray-triangle intersection test~\cite{moller1997} is performed along the
hull point's swept path.

The ray is defined as:
\begin{equation*}
\vec{O} = \vec{q} - \vec{v}_{\mathrm{rel}}^{\,\mathrm{local}} \cdot \Delta t, \qquad
\vec{D} = \vec{v}_{\mathrm{rel}}^{\,\mathrm{local}} \cdot \Delta t
\end{equation*}
where $\vec{O}$ is the hull point's estimated start-of-frame position in
mesh-local space, and $\vec{D}$ is the displacement during the frame.
The test is only performed when the velocity is approaching the surface
($v_{\hat{n}} = \vec{v}_{\mathrm{rel}} \cdot \hat{\vec{n}} < -0.01$).

The M\"oller-Trumbore algorithm computes:
\begin{equation*}
\vec{h} = \vec{D} \times \vec{e}_2, \quad
a = \vec{e}_1 \cdot \vec{h}
\end{equation*}
If $|a| < 10^{-8}$ the ray is parallel to the triangle and no intersection
is possible.  Otherwise the barycentric coordinates and parametric distance
are:
\begin{equation*}
f = 1/a, \quad
\vec{s} = \vec{O} - \vec{p}_0, \quad
u_r = f\,(\vec{s} \cdot \vec{h})
\end{equation*}
\begin{equation*}
\vec{q}_r = \vec{s} \times \vec{e}_1, \quad
v_r = f\,(\vec{D} \cdot \vec{q}_r), \quad
t_r = f\,(\vec{e}_2 \cdot \vec{q}_r)
\end{equation*}
A through-penetration is confirmed when $u_r \in [-0.05, 1.05]$,
$v_r \in [-0.05, 1.05]$, $u_r + v_r \le 1.10$, and $t_r \in [0, 1]$.
The barycentric tolerances match those used in the static test;
the parametric range $t_r \in [0,1]$ ensures the intersection occurred
during the current frame.


\subsection{Docking exclusion zones}
\label{ssec:dock_exclusion}
When two vessels are on a docking approach, collision detection must be
suppressed near the docking ports to allow the docking mechanism to
engage without triggering a collision response.

\subsubsection{Docking approach detection}
A docking approach is declared when any pair of dock ports
$(i \in A,\; j \in B)$ satisfies both:
\begin{itemize}
\item \textbf{Proximity}: $|\vec{p}_{d,A}^i - \vec{p}_{d,B}^j| < 5.0$\,m,
  where $\vec{p}_{d}$ is the dock port's reference position in global
  coordinates.
\item \textbf{Alignment}: $\vec{d}_{A}^i \cdot \vec{d}_{B}^j < -0.8$,
  where $\vec{d}$ is the dock port's approach direction.  A dot product
  below $-0.8$ indicates the ports face each other within $\sim\!37°$.
\end{itemize}

\subsubsection{Exclusion mechanisms}
When a docking approach is active, two exclusion mechanisms apply:

\begin{enumerate}
\item \textbf{Hull slice exclusion (convex paths).}
  Each \texttt{HullGroupSlice} carries an \texttt{isDockClearZone} flag.
  During a docking approach, slices marked as dock clear zones are
  skipped entirely in the convex-convex path.  In the convex-concave
  path, clear-zone slices are skipped for triangles whose vertices are
  within 3\,m of any dock port on the concave vessel.

\item \textbf{Per-point exclusion (concave fallback path).}
  In the point-to-triangle fallback, each hull point is transformed to
  planet-local space and checked against all dock port positions of the
  opposing vessel.  Points within 2\,m of any dock port are excluded
  from collision testing.
\end{enumerate}


\subsection{Collision response - rigid-body impulse solver}
\label{ssec:response}

\subsubsection{Contact velocity}
Once the deepest penetration point and surface normal $\hat{\vec{n}}$
are known, the contact-point velocity for each vessel is computed
including both linear and angular contributions:
\begin{equation}
\vec{v}_A^c = \vec{v}_A + \vec{r}_A \times \vec{\omega}_A, \qquad
\vec{v}_B^c = \vec{v}_B + \vec{r}_B \times \vec{\omega}_B
\label{eq:contact_vel}
\end{equation}
where $\vec{r}_A$ and $\vec{r}_B$ are the moment arms from each vessel's
centre of mass to the contact point.  Note that Orbiter uses a
left-handed coordinate convention where $\vec{v} = \vec{r} \times
\vec{\omega}$, which is reflected in the cross product ordering
throughout the impulse computation.

The relative contact velocity is
\begin{equation*}
\vec{v}_{\mathrm{rel}}^c = \vec{v}_A^c - \vec{v}_B^c
\end{equation*}
and the component along the collision normal is
\begin{equation*}
v_n = \vec{v}_{\mathrm{rel}}^c \cdot \hat{\vec{n}}
\end{equation*}
If $v_n > 0$ the bodies are already separating and no impulse is applied.
If $|v_n| > 500\,\text{m/s}$ a hypervelocity impact is declared and
both vessels are destroyed.

\subsubsection{Normal impulse magnitude}
The impulse magnitude $j$ along the collision normal is computed as
\begin{equation}
j = \frac{-(1+e)\,v_n}
     {\dfrac{1}{m_A} + \dfrac{1}{m_B}
      + \bigl[(\mat{I}_A^{-1}(\hat{\vec{n}} \times \vec{r}_A)) \times \vec{r}_A
        + (\mat{I}_B^{-1}(\hat{\vec{n}} \times \vec{r}_B)) \times \vec{r}_B\bigr]
        \cdot \hat{\vec{n}}}
\label{eq:impulse}
\end{equation}
where $e$ is the coefficient of restitution, $m_A$, $m_B$ are the vessel
masses, and $\mat{I}_A^{-1}$, $\mat{I}_B^{-1}$ are the inverse inertia
tensors expressed in the global frame.

For vessel-to-vessel collisions, $e = 1.0$ (perfectly elastic restitution).
Energy dissipation in vessel-to-vessel interactions is provided by
tangential friction (Section~\ref{sssec:friction}), angular damping
(Section~\ref{sssec:angdamp}), and the angular velocity cap
(Section~\ref{sssec:omegacap}) rather than through the restitution
coefficient.

The denominator in Eq.~\ref{eq:impulse} can be decomposed for clarity.
Define the \emph{angular effect} term:
\begin{equation}
\alpha = \bigl[(\mat{I}_A^{-1}\,\vec{c}_A) \times \vec{r}_A
         + (\mat{I}_B^{-1}\,\vec{c}_B) \times \vec{r}_B\bigr]
         \cdot \hat{\vec{n}}
\label{eq:angular_effect}
\end{equation}
where $\vec{c}_A = \hat{\vec{n}} \times \vec{r}_A$ and
$\vec{c}_B = \hat{\vec{n}} \times \vec{r}_B$.
Then Eq.~\ref{eq:impulse} simplifies to
\begin{equation}
j = \frac{-(1+e)\,v_n}{\dfrac{1}{m_A} + \dfrac{1}{m_B} + \alpha}
\label{eq:impulse_simple}
\end{equation}

The measured penetration depth is clamped to a maximum of 2.0\,m before
being used for positional correction, to prevent catastrophic
overcorrection when large interpenetrations occur due to low framerates or
high-speed impacts.

\subsubsection{Applying the normal impulse}
The velocity and angular velocity changes are
\begin{equation}
\begin{aligned}
\Delta\vec{v}_A &= +\frac{j\,\hat{\vec{n}}}{m_A}, &\qquad
\Delta\vec{v}_B &= -\frac{j\,\hat{\vec{n}}}{m_B} \\[4pt]
\Delta\vec{\omega}_A &= +\mat{I}_A^{-1}\,\vec{c}_A\;j, &\qquad
\Delta\vec{\omega}_B &= -\mat{I}_B^{-1}\,\vec{c}_B\;j
\end{aligned}
\label{eq:apply_impulse}
\end{equation}

\subsubsection{Tangential friction}
\label{sssec:friction}
The tangential component of the relative contact velocity is
\begin{equation*}
\vec{v}_t = \vec{v}_{\mathrm{rel}}^c - v_n\,\hat{\vec{n}}
\end{equation*}
A friction impulse is computed using the same impulse denominator structure
as the normal case but along the tangent direction
$\hat{\vec{t}} = \vec{v}_t / |\vec{v}_t|$:
\begin{equation}
j_t = \min\!\bigl(|\vec{v}_t|\,m_{\mathrm{eff}}^t,\;\mu\,j\bigr)
\label{eq:friction}
\end{equation}
where $\mu$ is the Coulomb friction coefficient (0.6 for vessel-to-vessel)
and $m_{\mathrm{eff}}^t$ is the effective mass in the tangential direction,
computed analogously to the normal case.
The friction impulse is capped by $\mu\,j$ (Coulomb's law) to prevent
the tangential correction from reversing the sliding direction.

For vessel-to-base collisions, wheel brakes are modelled by scaling $\mu$
from a base value of 0.8 up to 2.8 at full brake application:
\begin{equation*}
\mu = 0.8 + 2.0\,b
\end{equation*}
where $b \in [0,1]$ is the average of left and right wheel brake levels.

\subsubsection{Angular damping (vessel-to-vessel)}
\label{sssec:angdamp}
For vessel-to-vessel collisions, angular damping is applied only during
sustained/resting contact (closing speed $< 0.5\,\text{m/s}$), and only
damps the spin component around the contact normal to avoid killing
legitimate tumble from oblique impacts:
\begin{equation*}
\omega_{\mathrm{damp}} = (\vec{\omega}_{\mathrm{rel}} \cdot \hat{\vec{n}}) \cdot \hat{\vec{n}} \cdot \min(0.3, 3\,\Delta t)
\end{equation*}
This is subtracted from vessel~$A$'s angular velocity and added to vessel~$B$'s,
equalising only the ``rolling'' spin around the contact normal.

\subsubsection{Angular velocity cap}
\label{sssec:omegacap}
To prevent numerical explosions where single-point contact generates
unbounded angular acceleration, the angular velocity gain from any single
collision is capped:
\begin{equation*}
|\vec{\omega}_{\mathrm{post}}| \le |\vec{\omega}_{\mathrm{pre}}| + \omega_{\max}
\end{equation*}
where $\omega_{\max} = 10.0\,\text{rad/s}$.  If the post-impulse angular
velocity magnitude would exceed this limit, the angular velocity change
$\Delta\vec{\omega}$ is rescaled to exactly reach the cap while preserving
its direction.  This flat cap is intentionally generous to permit
legitimate angular momentum transfer from oblique impacts while guarding
against pure numerical instabilities.

\subsubsection{Positional correction}
After all impulses are applied, a direct positional correction
pushes the interpenetrating bodies apart:
\begin{equation}
\begin{aligned}
\Delta\vec{p}_A &= +\hat{\vec{n}}\,(d_{\mathrm{pen}} \cdot 1.01 + \delta_v) \cdot
                    \frac{m_B}{m_A + m_B} \\[4pt]
\Delta\vec{p}_B &= -\hat{\vec{n}}\,(d_{\mathrm{pen}} \cdot 1.01 + \delta_v) \cdot
                    \frac{m_A}{m_A + m_B}
\end{aligned}
\label{eq:pos_correction}
\end{equation}
where $d_{\mathrm{pen}}$ is the measured penetration depth and
$\delta_v = \min(v_{\mathrm{close}}\,\Delta t / 2,\; d_{\mathrm{pen}} / 2)$
is a velocity-proportional margin clamped to half the penetration depth
to prevent overshoot at low framerates.
The 1\% over-correction ($\times 1.01$) prevents persistent
frame-to-frame interpenetration in resting contact scenarios.


\subsection{Vessel-to-base collisions}
\label{ssec:vessel_base}
Vessel-to-base collisions follow the same point-to-triangle
algorithm described in Section~\ref{ssec:narrowphase_legacy}, using the
vessel's hull point cloud tested against base structure meshes.
The base is treated as an immovable object, so the impulse solver
simplifies: $m_B \to \infty$ and $\mat{I}_B^{-1} = \mat{0}$.
The normal impulse reduces to
\begin{equation*}
j_n = -(1 + e)\,v_n
\end{equation*}
with $e = 0.1$ for impacts exceeding 0.5\,m/s and $e = 0$ for softer
contacts to prevent resting-state jitter.

Angular damping is applied when a vessel is nearly at rest on a surface
($|v_n| < 0.5\,\text{m/s}$ and $d_{\mathrm{pen}} > 0.005\,\text{m}$):
the angular velocity relative to the planet's surface rotation is reduced
by 20\% per frame. This compensates for the single-contact-point
approximation and prevents tipping oscillations.
The 0.5\,m/s velocity threshold distinguishes ""resting" from "impacting"
states, as below this speed, contacts are treated as static friction
scenarios requiring stabilization; above it, normal restitution physics
apply. The 0.005\,m penetration requirement ensures damping only activates
when meaningful surface contact exists, preventing false activation during
grazing trajectories. The 20\% damping rate balances fast convergence to
rest (achieved in $\sim$10 frames) versus smooth, artifact-free motion
during slow rolling.


\subsection{Base configuration - the COLLISION marker}
\label{ssec:cfg_collision}
Base configuration files (\texttt{.cfg}) can optionally flag individual
mesh entries for collision detection using the \texttt{COLLISION} keyword.
If a mesh block does not contain this marker, no collision check is
performed against it (preserving legacy behaviour).

\begin{lstlisting}[language=OSFS]
HANGAR
  POS 0 0 0
  ROT 90
  SCALE 12 8 244
  TEX1 concrete1 0.4 0.4
  TEX3 concrete1 0.5 8
  COLLISION
END
\end{lstlisting}

The \texttt{COLLISION} keyword can appear anywhere within the mesh block
between its header and the closing \texttt{END} statement.
Multiple mesh entries in the same base configuration may independently
use or omit the marker.


\subsection{Diagnostic logging}
\label{ssec:debug}
The collision system includes a compile-time diagnostic mode activated
by defining \texttt{VESSEL\_COLLISION\_DEBUG} before compiling
\texttt{Vessel.cpp}.
When enabled, each vessel-to-vessel collision event writes detailed
diagnostic data to the Orbiter log file, including:
\begin{itemize}
\item Collision normal, penetration depth, and coefficient of restitution
\item Contact-point lever arms for both vessels
\item Impulse magnitude, effective masses, and angular effect terms
\item Pre- and post-impulse kinetic energy (linear + rotational)
\item Pre- and post-impulse momentum vectors
\item Energy gain/loss percentage and momentum conservation error
\item Friction impulse parameters (tangential velocity, friction coefficient)
\item Position correction magnitude and velocity margin
\end{itemize}
A \texttt{WARNING: ENERGY GAIN} message is emitted whenever the
post-impulse kinetic energy exceeds the pre-impulse value, indicating
a potential physics violation for further investigation.

\end{document}
